% Chapter 3: Jensen's Operator Inequality
\chapter{Jensen's Operator Inequality}

Jensen's operator inequality is the central engine of this formalization.
It generalizes the scalar Jensen inequality to the noncommutative setting: given an
operator convex function $f$ and a family of self-adjoint elements $a_i$ in an ordered
C$^*$-algebra, it asserts that $f$ of any ``weighted combination'' of the $a_i$ is
dominated by the corresponding weighted combination of the $f(a_i)$.
In this project the inequality is applied in Chapter~5 to deduce joint convexity and
concavity of the generalized perspective function, which in turn yields Lieb's concavity
theorem in Chapter~6.
We follow the unitary-averaging proof of Li and Wu~\cite{LiWu2012}.

\section{Main Theorem}

\begin{theorem}[Jensen's operator inequality]
  \label{thm:jensen}
  \lean{JensenOperator_convex_general}
  \leanok
  \uses{def:operator_convex}
  Let $\mathcal{B}$ be an ordered unital C$^*$-algebra, $I \subseteq \mathbb{R}$ (not necessarily an
  interval), and $f : \mathbb{R} \to \mathbb{R}$ a continuous operator convex function on
  $I$.
  Let $a_0, \ldots, a_n \in \mathcal{B}$ be self-adjoint elements with
  $\sigma(a_i) \subseteq I$ for each $i$, and let $b_0, \ldots, b_n \in \mathcal{B}$
  satisfy the unital condition
  \[
    \sum_{i=0}^{n} b_i^* b_i = 1.
  \]
  Then
  \[
    f\!\left(\sum_{i=0}^{n} b_i^* a_i\, b_i\right)
    \;\leq\;
    \sum_{i=0}^{n} b_i^*\, f(a_i)\, b_i.
  \]

  \begin{proof}
    \proves{thm:jensen}
    \leanok
    \uses{lem:conj_sum_spectrum, thm:fourier_avg, thm:cfc_unitary_conj, thm:cfc_diagonal, thm:diagonal_le_diagonal}
    We follow the unitary-averaging argument of Li and Wu.

    Write $m = n + 2$, set $a_{n+1} := a_n$, and form the diagonal block element
    $D = \operatorname{diag}(a_0, \ldots, a_{n+1}) \in M_m(\mathcal{B})$.
    By \ref{thm:cfc_diagonal} the continuous functional calculus commutes with block-diagonal
    structure, so $f(D) = \operatorname{diag}(f(a_0), \ldots, f(a_{n+1}))$.
    Because each $a_i$ has spectrum in $I$, the block element $D$ is self-adjoint with
    spectrum in $I$ by \ref{lem:conj_sum_spectrum}.

    Next, using the $b_i$ and the constraint $\sum_{i=0}^n b_i^* b_i = 1$, one constructs
    an explicit $(m \times m)$ matrix $U \in M_m(\mathcal{B})$ whose last column is
    $(b_0, \ldots, b_n, 0)^T$ and whose remaining columns are determined so that
    $U^* U = U U^* = 1$, i.e. $U$ is unitary.
    A direct computation using $\sum_i b_i^* b_i = 1$ shows that the bottom-right
    corner of $U^* D U$ equals $\sum_{i=0}^n b_i^* a_i b_i$.
    Since $U$ is unitary, $U^* D U$ is self-adjoint with spectrum equal to that of $D$,
    hence contained in $I$.

    Fix $\omega = e^{2\pi i/m}$, a primitive $m$-th root of unity, and for
    $k = 0, \ldots, m-1$ let
    $V_k = \operatorname{diag}(1, \omega^k, \omega^{2k}, \ldots, \omega^{(m-1)k}) \in M_m(\mathcal{B})$.
    Each $V_k$ is unitary (scalar multiples of the identity on each diagonal block).
    By \ref{thm:fourier_avg} we therefore have
    \[
      \frac{1}{m} \sum_{k=0}^{m-1} V_k^*\,(U^* D U)\,V_k
      \;=\; \operatorname{diag}\!\bigl((U^* D U)_{00},\, \ldots,\, (U^* D U)_{m-1,m-1}\bigr),
    \]
    whose bottom-right corner is still $\sum_{i=0}^n b_i^* a_i b_i$.

    Applying the same Fourier averaging to $U^* f(D) U$ and using \ref{thm:cfc_unitary_conj}
    (which gives $f(U^* D U) = U^* f(D) U$ for unitary $U$) together with \ref{thm:cfc_diagonal}
    yields
    \[
      \frac{1}{m} \sum_{k=0}^{m-1} V_k^*\,(U^* f(D) U)\,V_k\bigg|_{\mathrm{br}}
      \;=\; \sum_{i=0}^n b_i^*\, f(a_i)\, b_i.
    \]

    The Fourier average $\frac{1}{m}\sum_k V_k^*(U^* D U)V_k$ is a convex combination
    (uniform weights $\frac{1}{m}$) of elements $W_k^* D W_k$ with $W_k = U V_k$ unitary,
    each having spectrum in $I$.
    Operator convexity of $f$ on $M_m(\mathcal{B})$ (which follows from operator convexity
    on $\mathcal{B}$ since $f$ is assumed $\mathcal{B}$-convex of all orders) therefore gives
    \[
      f\!\left(\frac{1}{m}\sum_{k=0}^{m-1} V_k^*(U^* D U)V_k\right)
      \;\leq\;
      \frac{1}{m}\sum_{k=0}^{m-1} f\!\bigl(V_k^*(U^* D U)V_k\bigr)
      \;=\;
      \frac{1}{m}\sum_{k=0}^{m-1} V_k^*(U^* f(D) U)V_k,
    \]
    where the last equality uses \ref{thm:cfc_unitary_conj} applied to $U$ and each $V_k$.

    Applying \ref{thm:diagonal_le_diagonal} to the inequality above and substituting the identities above, specializing to 
    botom-right corner yields the inequality
    \[
      f\!\left(\sum_{i=0}^{n} b_i^* a_i b_i\right)
      \;\leq\;
      \sum_{i=0}^{n} b_i^*\, f(a_i)\, b_i.
    \]
  \end{proof}
\end{theorem}

\section{Sub-Unital Variant}

In applications it is often the case that the operators $b_i$ satisfy only a
sub-unital condition $\sum_i b_i^* b_i \leq 1$ rather than equality.
The following corollary handles this by padding the family with one extra element.

\begin{theorem}[Jensen's operator inequality, sub-unital case]
  \label{thm:jensen_sub}
  \lean{JensenOperator_convex_general_sub}
  \leanok
  \uses{def:operator_convex}
  Let $\mathcal{B}$, $I$, $f$, $a_i$, and $b_i$ be as in Theorem~\ref{thm:jensen},
  except that the $b_i$ now satisfy only
  \[
    \sum_{i=0}^{n} b_i^* b_i \leq 1.
  \]
  Assume in addition that $0 \in I$ and $f(0) \leq 0$.
  Then
  \[
    f\!\left(\sum_{i=0}^{n} b_i^* a_i\, b_i\right)
    \;\leq\;
    \sum_{i=0}^{n} b_i^*\, f(a_i)\, b_i.
  \]

  \begin{proof}
    \proves{thm:jensen_sub}
    \leanok
    \uses{thm:jensen}
    Set $b_{n+1} = \bigl(1 - \sum_{i=0}^{n} b_i^* b_i\bigr)^{1/2}$ and $a_{n+1} = 0$.
    The extended family satisfies $\sum_{i=0}^{n+1} b_i^* b_i = 1$ and
    $\sigma(a_{n+1}) = \{0\} \subseteq I$, so Theorem~\ref{thm:jensen} applies to give
    \[
      f\!\left(\sum_{i=0}^{n+1} b_i^* a_i\, b_i\right)
      \;\leq\;
      \sum_{i=0}^{n+1} b_i^*\, f(a_i)\, b_i.
    \]
    The left-hand side equals $f\!\bigl(\sum_{i=0}^{n} b_i^* a_i b_i\bigr)$ since
    $a_{n+1} = 0$.
    The last term on the right is $b_{n+1}^* f(0)\, b_{n+1} \leq 0$ by the assumption
    $f(0) \leq 0$, so dropping it only strengthens the inequality, yielding the
    claimed bound.
  \end{proof}
\end{theorem}
